\subsection{Craving-Übertragung und Tokenmechanismus}
Die inhaltlich relevante Selbstberichtsvariable der App ist der Craving-Wert. Er wird als ganzzahliger Wert zwischen 0 und 100 geführt und erst nach Abschluss der fünf Aufgaben eines Durchgangs übertragen. Die App sendet den Wert als JSON-Payload per HTTP-\texttt{PUT} an \texttt{submit.php}. Zusätzlich enthält der initiale Selbstbericht das lokale \texttt{app\_token}, den aktuell gültigen HMAC-Kettenwert \texttt{hash} und die App-Version. Die Endpunkt-URL wird über die Build-Konfiguration gesetzt: die Staging-Variante verwendet einen Test-Endpunkt, die Produktionsvariante \texttt{https://cuelens.each-and-every.de/submit}. Netzwerkanfragen werden in einem I/O-Coroutine-Kontext ausgeführt und blockieren daher nicht den UI-Thread.

Der Tokenmechanismus begrenzt die Zahl gültiger Selbstberichte, behandelt Wiederholungen nach Netzwerkfehlern idempotent und vermeidet eine dauerhafte Verbindung zwischen Registrierungsdaten und auswertbaren Selbstberichten. Bei der initialen Aktivierung sendet die App einmalig die E"~Mail-Adresse aus der Registrierung. Der Server prüft, ob diese Adresse in der Tabelle \texttt{register} vorhanden, bestätigt und noch nicht für ein App-Token verwendet wurde. Anschließend erzeugt der Server ein UUID-basiertes \texttt{app\_token} und den ersten HMAC-Kettenwert, speichert das App-Token aber nicht dauerhaft. Erst nachdem die App den Erhalt mit einem zweiten \texttt{PUT} bestätigt hat, wird in \texttt{register} lediglich ein technischer Zeitstempel gesetzt und der erste gültige Hash in \texttt{valid\_hashes} gespeichert.

Die Hash-Kette wird serverseitig aus dem App-Token \(t\) und einem nur serverseitig gespeicherten Schlüssel \(s\) berechnet:
\[
  h_1 = \mathrm{HMAC}_s(t), \qquad
  h_i = \mathrm{HMAC}_s(h_{i-1} \Vert t) \quad \text{für } i=2,\ldots,20.
\]
Der Server speichert immer nur den nächsten vorzulegenden Kettenwert in \texttt{valid\_hashes}. Aus dieser Tabelle allein sind weder E"~Mail-Adresse noch App-Token, Bedingung oder aktueller Studienfortschritt ersichtlich. Bei einem Selbstbericht berechnet der Server die Kette erneut, prüft den übermittelten Hash mit \texttt{hash\_equals} gegen alle 20 möglichen Kettenwerte und bricht die Prüfschleife nicht beim ersten Treffer ab. Aus dem gefundenen Index wird die Studienbedingung abgeleitet: Die Kettenindizes 1 bis 10 entsprechen Cue-Matching, die Indizes 11 bis 20 Cue-Labeling.

Die Craving-Übertragung folgt einem Drei-Wege-Handshake. Im ersten Schritt sendet die App \texttt{app\_token}, \texttt{hash}, \texttt{craving} und optional \texttt{app\_version}. Der Server prüft, ob der Hash in \texttt{valid\_hashes} existiert und zur berechneten Kette passt. Für die Kettenindizes 1 bis 19 wird der verbrauchte Hash gelöscht und der Selbstbericht zunächst nur temporär in \texttt{submission} gespeichert. Die Antwort enthält den nächsten Hash. Die App speichert diesen Wert zunächst als \texttt{pending\_confirmation\_hash} und sendet im zweiten Schritt einen Bestätigungs-\texttt{PUT}, der keine neuen Studiendaten enthält. Erst nach dieser Bestätigung übernimmt der Server den temporären Datensatz in \texttt{self\_reports}, speichert den bestätigten nächsten Hash als neuen gültigen Hash und löscht den Eintrag aus \texttt{submission}. Der Bestätigungsrequest liefert im Erfolgsfall und bei Wiederholung \texttt{HTTP 204}. Dadurch kann die App nach Verbindungsabbrüchen denselben Schritt wiederholen, ohne den Selbstbericht doppelt zu speichern.

Für den zwanzigsten Kettenwert wird kein weiterer Hash erzeugt. Der Server speichert den letzten Selbstbericht, erzeugt einen UUID-basierten Abrechnungscode und legt diesen in \texttt{compensation\_code} ab. Die App speichert den Code lokal, bestätigt ihn mit einem weiteren \texttt{PUT} und markiert die Studie erst nach erfolgreicher Bestätigung als abgeschlossen. Ein bereits bestätigter oder unbekannter Hash kann danach keinen zusätzlichen Selbstbericht erzeugen.

Serverseitig werden alle Requests als JSON-Objekt geparst; andere HTTP-Methoden als \texttt{PUT} führen zu \texttt{HTTP 405}. Der Craving-Wert wird als Ganzzahl validiert und auf den Bereich 0 bis 100 beschränkt. App-Token und Abrechnungscodes müssen dem UUID-Format entsprechen, HMAC-Werte müssen als 64-stellige Hex-Zeichenketten vorliegen. Für Datenbankzugriffe wird PDO mit deaktivierter Emulation vorbereiteter Statements verwendet. Zustandsänderungen, die Hashes verbrauchen, Registrierungen aktivieren oder temporäre Submissions finalisieren, werden innerhalb von Transaktionen und mit \texttt{FOR UPDATE} ausgeführt. Datenbankfehler werden nach außen nur als generische Fehlermeldung ausgegeben.

Die lokale Retry-Logik bildet den Serververtrag nach. Vor einem Selbstbericht wird zunächst ein \texttt{pending\_submission\_craving} persistiert. Geht die Serverantwort vor dem nächsten Hash verloren, wiederholt die App denselben initialen \texttt{PUT}. Liegt ein \texttt{pending\_confirmation\_hash} vor, wird dagegen nur noch der Bestätigungs-\texttt{PUT} wiederholt. Erst nach \texttt{HTTP 204} wird \texttt{current\_hash} ersetzt, der Zähler bestätigter Studiensituationen erhöht, die nächste Sperrzeit gesetzt und der ausstehende lokale Zustand gelöscht. Beim nächsten App-Start wird diese Wiederherstellung automatisch ausgeführt, bevor eine weitere Studiensituation begonnen werden kann.
